% ---------------------------------------------------------------------------
% DRAFT 7 — Dataset Construction + Experimental Setup.
% Base text reused from acl_source/acl_latex.tex per STRUCTURE.md's own
% "Keep"/"Keep, trim" verdict -- these sections were not identified as
% needing a rewrite. Fixes applied on top of the reused text:
%   - 192-vs-34-document claim corrected and explicitly reported (Phase 1.1)
%   - DeepSeek renamed to DeepSeek-R1-Distill-Llama-70B
%   - full model-identity table added, resolving the "documented in the
%     evaluation code" claim with the actual identifiers, verified in
%     scripts/evaluate.py and scripts/evaluate_new_models.py
%   - Gemma's exact tag (gemma3:4b) stated explicitly, with the provenance
%     note explaining why it needed re-confirming
% ---------------------------------------------------------------------------

\section{Dataset Construction}

Closing this gap requires building from Indian regulatory documents at a scale sufficient for
reliable model comparison, while preserving the domain expertise that interpreting them
demands---the constraint that shaped every design choice below.

\subsection{Source Document Collection}

We collected 192 regulatory documents spanning 1992--2026 from two official Indian government
sources: the Securities and Exchange Board of India (92 documents: circulars, master circulars,
regulations, orders; sebi.gov.in) and the Reserve Bank of India (100 documents: circulars, policy
statements, directions; rbi.org.in). Documents were downloaded using a custom Python scraping
pipeline and converted to clean text using pdfplumber, which handles the multi-column layouts and
embedded tables common in Indian regulatory PDFs particularly well. Figure~\ref{fig:pipeline}
(Appendix~\ref{app:examples}) illustrates the full dataset construction pipeline;
Appendix~\ref{app:sourcedocs} lists the regulatory categories covered.

The 406 annotated items are drawn from \textbf{34 of these 192 documents}---the collection was
sized for coverage of the regulatory categories in Appendix~\ref{app:sourcedocs}, while annotation
effort concentrated on documents dense enough in extractable QA content to be worth an expert
annotator's time. Section~\ref{app:provenance} reports the full provenance: item counts per
document (median 6, largest 63) and the document-level clustering used for the significance testing
in Section~\ref{sec:significance}, since items sharing a source document are not independent draws.

\subsection{Task Types}

IndiaFinBench defines four task types, each probing a distinct reasoning capability
(Table~\ref{tab:tasktypes}); Figure~\ref{fig:taskexamples} (Appendix~\ref{app:examples}) gives a
full worked example of each.

\begin{table}[t]
\centering\small
\begin{tabular}{lrp{9.7cm}}
\toprule
\textbf{Task} & \textbf{N} & \textbf{Definition} \\
\midrule
Regulatory Interpretation (REG) & 174 & Identify the correct rule, compliance threshold, or scope
of applicability from a regulatory passage---e.g., that a stock exchange must forward a registration
application ``not later than thirty days of receipt.'' \\
Numerical Reasoning (NUM) & 92 & Perform arithmetic over figures embedded in regulatory
text---e.g., a Small Finance Bank's maximum eligible dividend from its Tier 1 Capital Ratio and
adjusted profit after tax---requiring correct extraction \emph{and} execution. \\
Contradiction Detection (CON) & 62 & Given two passages from different regulatory instruments,
determine (Yes/No, one-sentence explanation) whether they contradict, probing supersession
tracking across frequently-amended circulars. \\
Temporal Reasoning (TMP) & 78 & Establish chronological ordering of regulatory events, identify
which rule version was in force at a given time, or compute elapsed time between milestones
referenced only by date. \\
\bottomrule
\end{tabular}
\caption{IndiaFinBench task types.}
\label{tab:tasktypes}
\end{table}

\subsection{Annotation Protocol}

All question-answer pairs were authored by the primary annotator, who has specialist experience with
Indian financial regulatory documents spanning SEBI's securities regulation framework, RBI's
prudential norms, and the amendment-chain structure linking hundreds of circulars over three
decades---expertise unavailable to crowdsourced pools, so expert single-annotator construction is
standard for domain-specific benchmarks \citep{hendrycks2021cuad, malik2021ildc,
islam2023financebench}, validated empirically by an independent human inter-annotator agreement
study (Section~\ref{sec:iaa}). Each item was individually reviewed for an unambiguous, unique answer
derivable from context alone, and consists of a context passage (80--500 words), a question, a
reference answer, and metadata fields (task type, difficulty, source document).

Answer formats are standardized by task type: extractive spans for regulatory interpretation and
temporal reasoning; calculated values with units for numerical reasoning (e.g., `Rs.~5,500 crore');
and `Yes' or `No' with a brief explanation for contradiction detection.

Difficulty levels were assigned by the number of reasoning steps required: easy items are
single-step extraction (160 items, 39.4\%), medium items require multi-clause reasoning or
calculation (182, 44.8\%), and hard items span multiple instruments or complex arithmetic (64,
15.8\%; Section~\ref{sec:difficulty} finds this label predicts empirical difficulty for only three
of the twelve evaluated models). The 406-item scale reflects a deliberate quality-over-quantity
decision, structurally incompatible with crowdsourcing (cf.\ FinanceBench's 150 expert-verified
items, Section~\ref{sec:related}).

\subsection{Model-Based Secondary Validation}

To confirm that items are unambiguously answerable from context, a secondary validation pass used
LLaMA-3.3-70B (via Groq API) as a secondary model-based quality check under a context-only, zero-shot
prompt (temperature = 0), accessed in isolation from the evaluation pipeline. Although LLaMA-3.3-70B
also appears in the main evaluation, the two uses are functionally distinct: this pass asks whether a
question is unambiguously answerable from its context passage, a different task from the
evaluation's open-ended QA, following established benchmark construction practice
\citep{islam2023financebench, hendrycks2021cuad}.

\begin{table}[t]
\centering\small
\begin{tabular}{lrrrr}
\toprule
 & \multicolumn{2}{c}{\textbf{Model-based}} & \multicolumn{2}{c}{\textbf{Human IAA}} \\
\textbf{Task} & \textbf{Items} & \textbf{Agr.} & \textbf{Items} & \textbf{Agr.} \\
\midrule
REG & 53 & 85.7\% & 63 & 92.1\% \\
NUM & 32 & 84.4\% & 44 & 81.8\% \\
CON & 30 & 96.7\%\textsuperscript{$\kappa$=.918} & 35 & 80.0\%\textsuperscript{$\kappa$=.712} \\
TMP & 35 & 77.1\% & 38 & 86.8\% \\
\midrule
\textbf{All} & \textbf{150} & \textbf{90.7\%} & \textbf{180} & \textbf{86.1\%} \\
\bottomrule
\end{tabular}
\caption{Annotation validation: model-based secondary pass (150-item subset) and human
inter-annotator agreement (180 items, three 60-item rounds, second annotator blind to reference
answers). Cohen's $\kappa$ applies only to the binary CON task.}
\label{tab:validation}
\end{table}

The secondary pass agreed with the primary annotation on 90.7\% of the sampled items
(Table~\ref{tab:validation}, left); items with genuine disagreement ($\sim$1.3\% of the initial set)
were removed.

\subsection{Human Inter-Annotator Agreement}
\label{sec:iaa}

To provide rigorous empirical validation of single-annotator annotation quality, we conducted a
human inter-annotator agreement (IAA) study in which a second independent human annotator completed
three annotation rounds of 60 items each (180 items total; 44.3\% of the benchmark), drawn
proportionally from all four task types and weighted toward medium and hard difficulty. The
annotator worked from context passages and questions alone, without access to the primary
annotator's reference answers at any stage.

Agreement was then computed between the primary annotator's reference answers and the second
annotator's responses, using the same four-stage scoring procedure applied to model predictions
(Section~\ref{sec:scoring}); results are in Table~\ref{tab:validation} (right).

The $\kappa$ = 0.712 for contradiction detection falls in the `substantial agreement' band by
\citet{landiskoch1977} conventions, and is comparable to human agreement rates reported for similar
binary contradiction detection tasks in legal NLP. Numerical reasoning (81.8\%) and contradiction
detection (80.0\%) are the two lowest task-level agreement rates, both driven by the same mechanism:
reference answers are often terse final values, while the second annotator's answers frequently embed
a full derivation or restate the question context around the same number, so the strict scorer's
exact and numeric-set matching stages can miss an answer that is substantively identical but longer.
We manually reviewed and independently re-verified the arithmetic or logic of all eight discordant
numerical-reasoning items across the three rounds and recorded the review, item by item, in
\texttt{annotation/iaa/num\_discordance\_review.csv} (released with the benchmark): every case's
computed values or final conclusions were equivalent; disagreements arose solely from intermediate-step
presence, currency symbol style, unit labelling, date format, or reasoning-path verbosity, never from a
numeric or substantive-conclusion discrepancy.

\section{Experimental Setup}

\subsection{Models}
\label{sec:models}

We evaluate twelve models spanning a wide range of sizes, providers, and access modes on the full
406-item benchmark: Gemini 2.5 Flash, Gemini 2.5 Pro, Qwen3-32B, LLaMA-3.3-70B, Llama 4 Scout 17B,
Kimi K2, GPT-OSS 120B, GPT-OSS 20B, DeepSeek-R1-Distill-Llama-70B, LLaMA-3-8B, Mistral-7B, and
Gemma 3 4B (the last three run locally via Ollama v0.6.5, 4-bit quantization, on an RTX 4060,
8~GB; exact providers, API identifiers, and completion budgets: Table~\ref{tab:models},
Appendix~\ref{app:models}). All models were evaluated zero-shot with no fine-tuning or prompt
adaptation, at temperature 0.0 --- but \textbf{not} under a shared completion budget: the harness set
per-provider generation limits (200--2048 tokens) for cost reasons unrelated to any experimental
design choice, and Gemini 2.5 Pro received no explicit cap at all. Budget heterogeneity is a
plausible confound for any format-sensitivity claim this paper makes, and we treat it as one rather
than assuming it away. Appendix~\ref{app:models} reports the exact provider, API model identifier or
Ollama tag, and completion budget used for every model, drawn directly from the evaluation harness
rather than restated from memory; Appendix~\ref{app:truncation} separately quantifies how much of the
logged prediction text this budget heterogeneity, and the harness's own logging limits, actually
truncate, and reports the paper's headline result restricted to the truncation-free item subset.
Three additional models --- Claude 3 Haiku (150 items), NVIDIA Nemotron 3 Super 120B (140 items),
and Qwen3-235B (120 items) --- were evaluated on partial subsets during earlier exploratory passes
and are excluded from every result in this paper because they were never run on the full 406-item
benchmark; their raw predictions remain in the public release for transparency, unused for any
reported number.
% TODO(Plan v3 Phase 2.5): once the matched-budget re-run lands, add a sentence here pointing to its
% results section and quantifying how much of the strict-score spread budget alone explains.

\subsection{Prompting Strategy}

Our evaluation follows the extractive QA paradigm established by SQuAD \citep{rajpurkar2016squad},
in the same spirit as DROP's discrete-reasoning-over-text framing \citep{dua2019drop}, with context
passages provided directly to the model under zero-shot, context-only constraints. This isolates
the model's ability to reason about the supplied regulatory text from its ability to recall
memorized facts; it does not by itself establish that the underlying regulatory documents are
absent from any model's training corpus, only that the model cannot use such memorized content
\emph{instead of} the supplied passage without the context-only instruction being violated. All
models received a system prompt establishing the context-only constraint (full text:
Appendix~\ref{app:prompt}). Task-specific user prompts provided appropriate formatting instructions
for each task type. For
contradiction detection, both passages were labelled explicitly as `Passage A / Passage B'. For
numerical reasoning, models were instructed to show calculation steps and include units---an
instruction in tension with the system prompt's ``no preamble'' constraint, which is one source of
the format-versus-content gap that Section~\ref{sec:regime} measures directly. All models received
the identical prompts described above at temperature 0.0; they did \textbf{not} share a completion
budget (Appendix~\ref{app:models} reports each model's exact value), which is a separate,
independently disclosed source of output-length variation from the prompting tension just described.

We evaluate exclusively under zero-shot conditions, as this most closely reflects practical
deployment where users query models without domain-specific priming, and because it eliminates
confounds from example selection strategy and ordering effects. Few-shot results for four top
models are reported in Appendix~\ref{app:fewshot}.

\subsection{Scoring}
\label{sec:scoring}

Answers were scored using a multi-stage matching procedure applied in sequence: (1) exact match
after case-normalization and punctuation stripping; (2) fuzzy token match using RapidFuzz
\texttt{token\_set\_ratio} $\geq$ 0.72, applied when exact match fails, with the 0.72 threshold
established by manual inspection of 20 borderline cases and validated against adjacent thresholds
(0.65 and 0.80); (3) numerical extraction match, where the item is scored correct if the set of
numbers extracted from both the reference and prediction are identical (handling currency symbols,
comma separators, and units); and (4) Yes/No match for contradiction detection, comparing the
leading word of the prediction to the reference. We refer to accuracy under this procedure as
\textbf{strict} accuracy throughout.

Section~\ref{sec:regime} evaluates the same predictions under two further regimes, built from a
cross-model judge (Section~\ref{sec:judge}) that reviews every prediction on REG, NUM, and TMP---not
only the ones strict scoring marks wrong: \textbf{judge-only}, in which the judge's verdict is final
in both directions, and \textbf{judge-augmented}, the asymmetric composite strict-OR-judge (a
strict-correct answer counts even if the judge disputes it; a strict-wrong answer counts if the judge
accepts it). Judge-only and strict are the paper's two primary regimes; judge-augmented is an
operational sensitivity regime, included to quantify what happens when a judge is applied only to
overturn strict failures and never to reject strict-correct ones, not because we consider it
authoritative. We avoid ``corrected,'' ``semantic,'' or unqualified ``judge-audited'' throughout:
the judge is itself an imperfect measurement (Section~\ref{sec:judge} reports its own error rate
against human adjudication), not a ground truth that strict scoring falls short of, and each of the
three regimes needs its own name because they can and do disagree.

DeepSeek-R1-Distill-Llama-70B was accessed via OpenRouter through an endpoint that returns final
answers without exposed reasoning traces; its stored predictions contain no chain-of-thought markup,
so we make no claim about the content of its internal deliberation, only about its final output
under its own completion budget (2048 tokens, the largest of any model in the study --- see
Appendix~\ref{app:models}), which rules out output-budget truncation as an explanation for its
strict-scoring behaviour specifically.
